\renewcommand{\refname}{참고문헌}
\begin{thebibliography}{131}\small
\bibitem{r1} ITU-R, M.2160: Framework and overall objectives of the future development of IMT for 2030 and beyond, \url{https://www.itu.int/rec/R-REC-M.2160/en}, [Online] (2023).
\bibitem{r2} 3GPP, Service requirements for the 5G system, TS 22.261 V19.5.0., Release 19 (2023).
\bibitem{r3} 3GPP, Study on physical layer enhancements for NR ultra-reliable and low latency case (URLLC), TR 38.824 V16.0.0., Release 16 (2019).
\bibitem{r4} X. You, Y. Huang, S. Liu, D. Wang, J. Ma, C. Zhang, H. Zhan, C. Zhang, J. Zhang, Z. Liu, et al., Toward 6G TK$\mu$ extreme connectivity: Architecture, key technologies and experiments, IEEE Wireless Commun. 30 (3) (2023) 86--95.
\bibitem{r5} A. Pourkabirian, M. S. Kordafshari, A. Jindal, M. H. Anisi, A vision of 6G URLLC: Physical-layer technologies and enablers, IEEE Commun. Standards Mag. 8 (2) (2024) 20--27.
\bibitem{r6} Y. Huang, Challenges and opportunities of sub-6 GHz integrated sensing and communications for 5G-advanced and beyond, Chinese J. Electron. 33 (2) (2024) 323--325.
\bibitem{r7} Y. Li, J. Zhao, J. Liao, F. Hu, Cellular v2x-based integrated sensing and communication system: Feasibility and performance analysis, Chinese J. Electron. 33 (4) (2024) 1104--1116.
\bibitem{r8} C. E. Shannon, A mathematical theory of communication, Bell Syst. Tech. J. 27 (3) (1948) 379--423.
\bibitem{r9} I. Kontoyiannis, S. Verdú, Optimal lossless data compression: Non-asymptotics and asymptotics, IEEE Trans. Inform. Theory 60 (2) (2013) 777--795.
\bibitem{r10} V. Kostina, S. Verdú, Fixed-length lossy compression in the finite blocklength regime, IEEE Trans. Inform. Theory 58 (6) (2012) 3309--3338.
\bibitem{r11} Y. Polyanskiy, H. V. Poor, S. Verdú, Channel coding rate in the finite blocklength regime, IEEE Trans. Inform. Theory 56 (5) (2010) 2307--2359.
\bibitem{r12} V. Kostina, S. Verdú, Lossy joint source-channel coding in the finite blocklength regime, IEEE Trans. Inform. Theory 59 (5) (2013) 2545--2575.
\bibitem{r13} L. Zhou, M. Motani, et al., Finite blocklength lossy source coding for discrete memoryless sources, Foundations and Trends\textregistered{} in Communications and Information Theory 20 (3) (2023) 157--389.
\bibitem{r14} V. Kostina, Lossy data compression: nonasymptotic fundamental limits, Ph.D.\ thesis, Princeton University (2013).
\bibitem{r15} S. Verdú, Teaching it, in: Proc. IEEE Int. Symp. Inf. Theory (ISIT), IEEE, 2007, pp. ix--ix.
\bibitem{r16} I. Kontoyiannis, Pointwise redundancy in lossy data compression and universal lossy data compression, IEEE Trans. Inform. Theory 46 (1) (2000) 136--152.
\bibitem{r17} K. Marton, Error exponent for source coding with a fidelity criterion, IEEE Trans. Inform. Theory 20 (2) (1974) 197--199.
\bibitem{r18} T. S. Han, The reliability functions of the general source with fixed-length coding, IEEE Trans. Inform. Theory 46 (6) (2006) 2117--2132.
\bibitem{r19} K. Iriyama, Probability of error for the fixed-length source coding of general sources, IEEE Trans. Inform. Theory 47 (4) (2001) 1537--1543.
\bibitem{r20} E.-h. Yang, Z. Zhang, On the redundancy of lossy source coding with abstract alphabets, IEEE Trans. Inform. Theory 45 (4) (1999) 1092--1110.
\bibitem{r21} R. Gray, Information rates of autoregressive processes, IEEE Trans. Inform. Theory 16 (4) (1970) 412--421.
\bibitem{r22} P. Tian, V. Kostina, The dispersion of the Gauss--Markov source, IEEE Trans. Inform. Theory 65 (10) (2019) 6355--6384.
\bibitem{r23} P. Tian, V. Kostina, Nonstationary Gauss-Markov processes: Parameter estimation and dispersion, IEEE Trans. Inform. Theory 67 (4) (2021) 2426--2449.
\bibitem{r24} A. Lapidoth, On the role of mismatch in rate distortion theory, IEEE Trans. Inform. Theory 43 (1) (1997) 38--47.
\bibitem{r25} L. Zhou, V. Y. Tan, M. Motani, Refined asymptotics for rate-distortion using Gaussian codebooks for arbitrary sources, IEEE Trans. Inform. Theory 65 (5) (2018) 3145--3159.
\bibitem{r26} A. H. Kaspi, Rate-distortion function when side-information may be present at the decoder, IEEE Trans. Inform. Theory 40 (6) (2002) 2031--2034.
\bibitem{r27} L. Zhou, M. Motani, Non-asymptotic converse bounds and refined asymptotics for two source coding problems, IEEE Trans. Inform. Theory 65 (10) (2019) 6414--6440.
\bibitem{r28} A. No, A. Ingber, T. Weissman, Strong successive refinability and rate-distortion-complexity tradeoff, IEEE Trans. Inform. Theory 62 (6) (2016) 3618--3635.
\bibitem{r29} L. Zhou, V. Y. Tan, M. Motani, Second-order and moderate deviations asymptotics for successive refinement, IEEE Trans. Inform. Theory 63 (5) (2017) 2896--2921.
\bibitem{r30} B. Rimoldi, Successive refinement of information: Characterization of the achievable rates, IEEE Trans. Inform. Theory 40 (1) (1994) 253--259.
\bibitem{r31} M. Effros, Distortion-rate bounds for fixed-and variable-rate multiresolution source codes, IEEE Trans. Inform. Theory 45 (6) (1999) 1887--1910.
\bibitem{r32} W. Yang, A. Collins, G. Durisi, Y. Polyanskiy, H. V. Poor, Beta--beta bounds: Finite-blocklength analog of the golden formula, IEEE Trans. Inform. Theory 64 (9) (2018) 6236--6256.
\bibitem{r33} Y. Polyanskiy, H. V. Poor, S. Verdu, Feedback in the non-asymptotic regime, IEEE Trans. Inform. Theory 57 (8) (2011) 4903--4925.
\bibitem{r34} R. C. Yavas, V. Y. Tan, Variable-length feedback codes over known and unknown channels with non-vanishing error probabilities, IEEE Trans. Inform. Theory 71 (5) (2025) 3271--3286.
\bibitem{r35} Y. Polyanskiy, S. Verdú, Arimoto channel coding converse and Rényi divergence, in: Proc. 48th Allerton Conf. Commun., IEEE, 2010, pp. 1327--1333.
\bibitem{r36} W. Yang, G. Durisi, T. Koch, Y. Polyanskiy, Quasi-static multiple-antenna fading channels at finite blocklength, IEEE Trans. Inform. Theory 60 (7) (2014) 4232--4265.
\bibitem{r37} Y. Polyanskiy, S. Verdú, Scalar coherent fading channel: Dispersion analysis, in: Proc. IEEE Int. Symp. Inf. Theory (ISIT), IEEE, 2011, pp. 2959--2963.
\bibitem{r38} A. Collins, Y. Polyanskiy, Coherent multiple-antenna block-fading channels at finite blocklength, IEEE Trans. Inform. Theory 65 (1) (2019) 380--405.
\bibitem{r39} J. Hoydis, R. Couillet, P. Piantanida, The second-order coding rate of the MIMO quasi-static Rayleigh fading channel, IEEE Trans. Inform. Theory 61 (12) (2015) 6591--6622.
\bibitem{r40} X. You, B. Sheng, Y. Huang, W. Xu, C. Zhang, D. Wang, P. Zhu, C. Ji, Closed-form approximation for performance bound of finite blocklength massive MIMO transmission, IEEE Trans. Commun. 71 (12) (2023) 6939--6951.
\bibitem{r41} Y. Polyanskiy, S. Verdú, Channel dispersion and moderate deviations limits for memoryless channels, in: Proc. Allerton Conf. Commun., Contr., Comput., IEEE, 2010, pp. 1334--1339.
\bibitem{r42} Y. Altuǧ, A. B. Wagner, Moderate deviations in channel coding, IEEE Trans. Inform. Theory 60 (8) (2014) 4417--4426.
\bibitem{r43} E. MolavianJazi, J. N. Laneman, A second-order achievable rate region for Gaussian multi-access channels via a central limit theorem for functions, IEEE Trans. Inform. Theory 61 (12) (2015) 6719--6733.
\bibitem{r44} Y.-W. Huang, P. Moulin, Finite blocklength coding for multiple access channels, in: Proc. IEEE Int. Symp. Inf. Theory (ISIT), IEEE, 2012, pp. 831--835.
\bibitem{r45} E. Haim, Y. Kochman, U. Erez, A note on the dispersion of network problems, in: Proc. IEEE Conv. Elect. Electron. Eng. Israel (IEEEI), IEEE, 2012, pp. 1--9.
\bibitem{r46} J. Scarlett, A. Martinez, A. G. i Fàbregas, Second-order rate region of constant-composition codes for the multiple-access channel, IEEE Trans. Inform. Theory 61 (1) (2015) 157--172.
\bibitem{r47} J. Scarlett, V. Y. Tan, Second-order asymptotics for the gaussian mac with degraded message sets, IEEE Trans. Inform. Theory 61 (12) (2015) 6700--6718.
\bibitem{r48} Y. Polyanskiy, A perspective on massive random-access, in: Proc. IEEE Int. Symp. Inf. Theory (ISIT), IEEE, 2017, pp. 2523--2527.
\bibitem{r49} S. S. Kowshik, K. Andreev, A. Frolov, Y. Polyanskiy, Energy efficient coded random access for the wireless uplink, IEEE Trans. Commun. 68 (8) (2020) 4694--4708.
\bibitem{r50} Y. Han, Z. Xiong, A. Høst-Madsen, Finite blocklength analysis of energy-delay tradeoff in uncoordinated mac, IEEE Trans. Commun. (2025).
\bibitem{r51} K.-H. Ngo, A. Lancho, G. Durisi, A. G. i Amat, Unsourced multiple access with random user activity, IEEE Trans. Inform. Theory 69 (7) (2023) 4537--4558.
\bibitem{r52} J. Gao, Y. Wu, S. Shao, W. Yang, H. V. Poor, Energy efficiency of massive random access in MIMO quasi-static Rayleigh fading channels with finite blocklength, IEEE Trans. Inform. Theory 69 (3) (2023) 1618--1657.
\bibitem{r53} J. Gao, Y. Wu, T. Li, W. Zhang, Energy efficiency of MIMO massive unsourced random access with finite blocklength, IEEE Wireless Commun. Lett. 12 (4) (2023) 743--747.
\bibitem{r54} J. Gao, Y. Wu, G. Caire, W. Yang, H. V. Poor, W. Zhang, Unsourced random access in MIMO quasi-static Rayleigh fading channels: Finite blocklength and scaling law analyses, IEEE Trans. Inform. Theory, early access (2025).
\bibitem{r55} I. Csiszár, On the error exponent of source-channel transmission with a distortion threshold, IEEE Trans. Inform. Theory 28 (6) (1982) 823--828.
\bibitem{r56} R. Pilc, The transmission distortion of a source as a function of the encoding block length, Bell Syst. Tech. J. 47 (6) (1968) 827--885.
\bibitem{r57} A. D. Wyner, On the transmission of correlated Gaussian data over a noisy channel with finite encoding block length, Inf. Contr. 20 (3) (1972) 193--215.
\bibitem{r58} I. Csiszár, J. Korner, Graph decomposition: A new key to coding theorems, IEEE Trans. Inform. Theory 27 (1) (1981) 5--12.
\bibitem{r59} D. Wang, A. Ingber, Y. Kochman, The dispersion of joint source-channel coding, in: Proc. Allerton Conf. Commun., Control, Comput., IEEE, 2011, pp. 180--187.
\bibitem{r60} D. Gündüz, M. A. Wigger, T.-Y. Tung, P. Zhang, Y. Xiao, Joint source--channel coding: Fundamentals and recent progress in practical designs, Proceedings of the IEEE (2024).
\bibitem{r61} E. Bourtsoulatze, D. B. Kurka, D. Gündüz, Deep joint source-channel coding for wireless image transmission, IEEE Tran. Cogn. Netw. 5 (3) (2019) 567--579.
\bibitem{r62} D. B. Kurka, D. Gündüz, DeepJSCC-f: Deep joint source-channel coding of images with feedback, IEEE J. Sel. Areas Inf. Theory 1 (1) (2020) 178--193.
\bibitem{r63} G. Wade, Signal coding and processing, Cambridge university press, 1994.
\bibitem{r64} S. Verdú, Teaching lossless data compression, IEEE Inform. Theory Society Newsletter 61 (1) (2011) 18--19.
\bibitem{r65} Y. Polyanskiy, Channel coding: Non-asymptotic fundamental limits, Ph.D.\ thesis, Princeton University, Princeton, NJ, USA (2010).
\bibitem{r66} T. M. Cover, J. A. Thomas, Elements of information theory, 2nd ed. New Jersey: Wiley, 2006.
\bibitem{r67} I. Kontoyiannis, Second-order noiseless source coding theorems, IEEE Trans. Inform. Theory 43 (4) (1997) 1339--1341.
\bibitem{r68} N. Merhav, Universal coding with minimum probability of codeword length overflow, IEEE Trans. Inform. Theory 37 (3) (1991) 556--563.
\bibitem{r69} T. A. Welch, A technique for high-performance data compression, Computer 17 (06) (1984) 8--19.
\bibitem{r70} I. Csiszár, J. Körner, Information theory: coding theorems for discrete memoryless systems, Cambridge University Press, 2011.
\bibitem{r71} W. Szpankowski, S. Verdú, Minimum expected length of fixed-to-variable lossless compression without prefix constraints, IEEE Trans. Inform. Theory 57 (7) (2011) 4017--4025.
\bibitem{r72} A. A. Yushkevich, On limit theorems connected with the concept of entropy of Markov chains, Uspekhi Matematicheskikh Nauk 8 (5) (1953) 177--180.
\bibitem{r73} V. Strassen, Asymptotische abschätzugen in Shannon’s informationstheorie, in: Proc. Trans. 3rd Prague Conf. Inf. Theory, 1962, pp. 689--723.
\bibitem{r74} V. V. Petrov, Limit Theorems of Probability Theory: Sequences of Independent Random Variables, Oxford, U.K.: Oxford Science, 1995.
\bibitem{r75} C. E. Shannon, et al., Coding theorems for a discrete source with a fidelity criterion, IRE Nat. Conv. Rec 4 (142-163) (1959) 1.
\bibitem{r76} J. Körner, et al., Coding of an information source having ambiguous alphabet and the entropy of graphs, in: 6th Prague conference on Information Theory, Academia, Prague, 1971, pp. 411--425.
\bibitem{r77} J. C. Kieffer, Strong converses in source coding relative to a fidelity criterion, IEEE Trans. Inform. Theory 37 (2) (1991) 257--262.
\bibitem{r78} S. Verdú, ELE528: Information theory lecture notes, Princeton University (2009).
\bibitem{r79} T. J. Goblick, Coding for a discrete information source with a distortion measure, Ph.D.\ thesis, Massachusetts Institute of Technology (1963).
\bibitem{r80} J. T. Pinkston, Encoding independent sample information sources. (1967).
\bibitem{r81} D. Sakrison, A geometric treatment of the source encoding of a Gaussian random variable, IEEE Trans. Inform. Theory 14 (3) (1968) 481--486.
\bibitem{r82} D. Varodayan, A. Aaron, B. Girod, Exploiting spatial correlation in pixel-domain distributed image compression, in: Proc. Picture Coding Symposium, Beijing, China, 2006.
\bibitem{r83} N. D. Memon, K. Sayood, Lossless compression of video sequences, IEEE Trans. Commun. 44 (10) (1996) 1340--1345.
\bibitem{r84} A. Kolmogorov, On the shannon theory of information transmission in the case of continuous signals, IRE Transactions on Information Theory 2 (4) (1956) 102--108.
\bibitem{r85} T. Berger, Information rates of Wiener processes, IEEE Trans. Inform. Theory 16 (2) (1970) 134--139.
\bibitem{r86} L. D. Davisson, Rate-distortion theory and application, Proceedings of the IEEE 60 (7) (1972) 800--808.
\bibitem{r87} I. Dumer, M. S. Pinsker, V. V. Prelov, On coverings of ellipsoids in Euclidean spaces, IEEE Trans. Inform. Theory 50 (10) (2004) 2348--2356.
\bibitem{r88} W. H. Equitz, T. M. Cover, Successive refinement of information, IEEE Trans. Inform. Theory 37 (2) (1991) 269--275.
\bibitem{r89} V. Kostina, S. Verdú, A new converse in rate-distortion theory, in: 2012 46th Annual Conference on Information Sciences and Systems (CISS), IEEE, 2012, pp. 1--6.
\bibitem{r90} S.-Q. Le, V. Y. Tan, M. Motani, Second-order coding rates for conditional rate-distortion, arXiv preprint arXiv:1410.2687 (2014).
\bibitem{r91} A. Kanlis, P. Narayan, Error exponents for successive refinement by partitioning, IEEE Trans. Inform. Theory 42 (1) (1996) 275--282.
\bibitem{r92} E. Tuncel, K. Rose, Error exponents in scalable source coding, IEEE Trans. Inform. Theory 49 (1) (2003) 289--296.
\bibitem{r93} V. Y. Tan, O. Kosut, On the dispersions of three network information theory problems, IEEE Trans. Inform. Theory 60 (2) (2013) 881--903.
\bibitem{r94} A. Feinstein, A new basic theorem of information theory, IRE Trans. Inform. Theory 4 (4) (1954) 2--22.
\bibitem{r95} C. E. Shannon, Certain results in the coding theory for noisy channels, Inf. Contr. 1 (1957) 6--25.
\bibitem{r96} R. G. Gallager, A simple derivation of the coding theorem and some applications, IEEE Trans. Inform. Theory 11 (1) (1965) 3--18.
\bibitem{r97} J. Scarlett, Reliable communication under mismatched decoding, Ph.D.\ thesis, University of Cambridge, Trinity Hall, Cambridge, UK (2014).
\bibitem{r98} W. Yang, R. F. Schaefer, H. V. Poor, Wiretap channels: Nonasymptotic fundamental limits, IEEE Trans. Inform. Theory 65 (7) (2019) 4069--4093.
\bibitem{r99} C. E. Shannon, Probability of error for optimal codes in a Gaussian channel, Bell Syst. Tech. J. 38 (3) (1959) 611--656.
\bibitem{r100} R. G. Gallager, Information theory and reliable communication, Wiley, New York, 1968.
\bibitem{r101} J. Wolfowitz, The coding of messages subject to chance errors, Illinois J. Math. 1 (1957) 591--606.
\bibitem{r102} C. E. Shannon, R. G. Gallager, E. R. Berlekamp, Lower bounds to error probability for coding on discrete memoryless channels I, Inf. Contr. 10 (1967) 65--103.
\bibitem{r103} S. Verdú, T. S. Han, A general formula for channel capacity, IEEE Trans. Inform. Theory 40 (4) (1994) 1147--1157.
\bibitem{r104} L. Shen, Y. Wu, Y. Xu, X. You, X. Gao, W. Zhang, GLDPC-PC codes: Channel coding towards 6G communications, accepted by IEEE Commun. Mag., arXiv preprint arXiv:2404.14828 (2025).
\bibitem{r105} P. Yuan, K. R. Duffy, M. Médard, Soft-output successive cancellation list decoding, IEEE Trans. Inform. Theory 71 (2) (2025) 1007--1017.
\bibitem{r106} L. Shen, Y. Wu, Y. Xu, X. You, X. Gao, W. Zhang, Soft-output fast successive-cancellation list decoder for polar codes, in: Proc. IEEE Wireless Commun. Netw. Conf. (WCNC), IEEE, 2025, pp. 1--6.
\bibitem{r107} W. Yang, G. Durisi, T. Koch, Y. Polyanskiy, Diversity versus channel knowledge at finite block-length, in: Proc. IEEE Inf. Theory Workshop (ITW), IEEE, 2012, pp. 572--576.
\bibitem{r108} F. Ye, X. You, J. Li, C. Zhang, C. Ji, Explicit performance bound of finite blocklength coded MIMO: Time-domain versus spatiotemporal channel coding, arXiv preprint arXiv:2406.13922v2 (2024).
\bibitem{r109} X. You, 6G extreme connectivity via exploring spatiotemporal exchangeability, Sci. China Inf. Sci. 66 (3) (2023) 130306.
\bibitem{r110} R. C. Yavas, V. Kostina, M. Effros, Third-order analysis of channel coding in the small-to-moderate deviations regime, IEEE Trans. Inform. Theory (2024).
\bibitem{r111} A. D. Wyner, Recent results in the shannon theory, IEEE Trans. Inform. Theory 20 (1) (1974) 2--10.
\bibitem{r112} Y. Wu, X. Gao, S. Zhou, W. Yang, Y. Polyanskiy, G. Caire, Massive access for future wireless communication systems, IEEE Wireless Commun. 27 (4) (2020) 148--156.
\bibitem{r113} L. Liu, W. Yu, Massive connectivity with massive MIMO---Part I: Device activity detection and channel estimation, IEEE Trans. Signal Process. 66 (11) (2018) 2933--2946.
\bibitem{r114} L. Liu, E. G. Larsson, W. Yu, P. Popovski, C. Stefanovic, E. De Carvalho, Sparse signal processing for grant-free massive connectivity: A future paradigm for random access protocols in the Internet of Things, IEEE Signal Process. Mag. 35 (5) (2018) 88--99.
\bibitem{r115} X. Chen, T.-Y. Chen, D. Guo, Capacity of Gaussian many-access channels, IEEE Trans. Inform. Theory 63 (6) (2017) 3516--3539.
\bibitem{r116} I. Zadik, Y. Polyanskiy, C. Thrampoulidis, Improved bounds on Gaussian mac and sparse regression via Gaussian inequalities, in: Proc. IEEE Int. Symp. Inf. Theory (ISIT), IEEE, 2019, pp. 430--434.
\bibitem{r117} S. S. Kowshik, Y. Polyanskiy, Fundamental limits of many-user mac with finite payloads and fading, IEEE Trans. Inform. Theory 67 (9) (2021) 5853--5884.
\bibitem{r118} J. Gao, Y. Wu, W. Zhang, Energy-efficiency of massive random access with individual codebook, in: Proc. IEEE Global Commun. (GLOBECOM), IEEE, 2020, pp. 1--6.
\bibitem{r119} A. Fengler, S. Haghighatshoar, P. Jung, G. Caire, Non-bayesian activity detection, large-scale fading coefficient estimation, and unsourced random access with a massive MIMO receiver, IEEE Trans. Inform. Theory 67 (5) (2021) 2925--2951.
\bibitem{r120} M. J. Ahmadi, M. Kazemi, T. M. Duman, Unsourced random access using multiple stages of orthogonal pilots: MIMO and single-antenna structures, IEEE Trans. Wireless Commun. 23 (2) (2024) 1343--1355.
\bibitem{r121} A. Fengler, O. Musa, P. Jung, G. Caire, Pilot-based unsourced random access with a massive MIMO receiver, interference cancellation, and power control, IEEE J. Sel. Areas Commun. 40 (5) (2022) 1522--1534.
\bibitem{r122} M. Gkagkos, K. R. Narayanan, J.-F. Chamberland, C. N. Georghiades, Fasura: A scheme for quasi-static fading unsourced random access channels, IEEE Trans. Commun. 71 (11) (2023) 6391--6401.
\bibitem{r123} A. Decurninge, I. Land, M. Guillaud, Tensor-based modulation for unsourced massive random access, IEEE Wireless Commun. Lett. 10 (3) (2021) 552--556.
\bibitem{r124} R. J. Pilc, Coding theorems for discrete source-channel pairs, Ph.D.\ thesis, Massachusetts Institute of Technology (1967).
\bibitem{r125} A. Wyner, Communication of analog data from a Gaussian source over a noisy channel, Bell Syst. Tech. J. 47 (5) (1968) 801--812.
\bibitem{r126} A. T. Campo, G. Vazquez-Vilar, A. G. i Fabregas, A. Martinez, Random-coding joint source-channel bounds, in: Proc. IEEE Int. Symp. Inf. Theory (ISIT), IEEE, 2011, pp. 899--902.
\bibitem{r127} X. Wang, C. Li, Z. Sun, L. Hui, Review of GAN-based research on Chinese character font generation, Chinese J. Electron. 33 (3) (2024) 584--600.
\bibitem{r128} S. Yue, Y. Deng, G. Wang, J. Ren, Y. Zhang, Federated offline reinforcement learning with proximal policy evaluation, Chinese J. Electron. 33 (6) (2024) 1360--1372.
\bibitem{r129} L. Ozarow, On a source-coding problem with two channels and three receivers, Bell Syst. Tech. J. 59 (10) (1980) 1909--1921.
\bibitem{r130} P. Zhang, G. Shi, S. Cui, Z. Zhang, K. Niu, Y. Xiao, Z. Qin, J. Dai, S. Shao, G. Deniz, G. Eleonora, Semantic communications: Theories, technologies and applications, China Commun. 21 (7) (2024) iii--vii.
\bibitem{r131} H. Tang, H. Zhu, H. Wei, H. Zheng, X. Mao, M. Lu, J. Guo, Representation of semantic word embeddings based on SLDA and Word2vec model, Chinese J. Electron. 32 (3) (2023) 647--654.
\end{thebibliography}
